\documentclass[11pt]{article}
\usepackage[a4paper,margin=2.5cm]{geometry}

\setlength{\parindent}{0pt}

\begin{document}

\begin{center}
{\Large\bfseries Gradient Decoupling in Adaptive Loss Balancing for Physics-Informed Neural Network Inverse Problems: A Provable Failure Mode and Design Principle}
\end{center}

\vspace{0.8cm}

\textbf{Abstract}

\vspace{0.3cm}

Physics-informed neural networks (PINNs) embed physical laws as soft constraints during training, but their reliability is limited by multi-task gradient imbalance among the PDE, data, and boundary-condition losses. We identify and formally prove a failure mode of gradient-based multi-task loss balancing in PINN inverse problems. When a gradient-norm balancer (e.g., ReLoBRaLo) uses the last-layer weights as its anchor, any loss term decoupled from that anchor---such as a parameter prior---has exactly zero anchor gradient, so its balancing weight decays geometrically to zero and the inverse problem is re-ill-posed. Full-convergence experiments (30,000 epochs) on a shallow-water inverse problem quantify this failure: balancing the prior drives a physical-parameter error from 0.047\% to 37.2\%, while fixing the prior weight preserves sub-percent recovery. We further show, across three PDE benchmarks (shallow water, Burgers, Navier--Stokes), that causal time-domain weighting does not transfer to Adam-only optimization. These results yield a design principle---\emph{balance conflicting losses, never parameter priors}---with reproducible guidance for PINN-based inverse problems.

\end{document}
